\chapter{Conclusion and Future Work}
\label{ch:conclusion}

\section{Answer to the Main Research Question}
The study began with two questions about weld quality. Could a camera reveal changes that machine signals missed? Would the combined prediction remain reliable if either input deteriorated? The architecture comparison extended to gear monitoring and robotic assembly. By the time the Intel Welding data were re-examined, development had already drawn on the proposed test recordings. Independent welding validation would require a new set of welds. The five studies evaluated related questions against the outcomes available in each dataset (Appendix~\ref{app:chronology}).

In PHM, post-run photographs supplied a continuous measure of candidate damage area on gear teeth. The photographs defined the target for sensor regression. Across held-out experiments, combining three sensor groups gave a mean absolute error of 0.824 percentage points, compared with 0.804 for predicting the training mean (Table~\ref{tab:phm-loeo}). The paired interval allowed either model to have lower error. Even an accurate prediction of this proxy would still require independent measurements to establish how closely the photographed area tracked physical damage.

REASSEMBLE made a direct fusion comparison possible. Video and machine signals covered the same completed action, each action had a recorded outcome, and entire recordings could be held out for assessment. Across 4,530 actions from 148 recordings, logistic stacking achieved failure average precision (AP) of .447, compared with .429 for its sensor component and .446 for equal probability averaging (Tables~\ref{tab:reassemble-unimodal} and~\ref{tab:reassemble-fusion}). The recording-level bootstrap supported a gain over sensors for these fitted models. The difference between stacking and averaging remained uncertain. Because the sensor predictor lacked a matched second calibration stage, the comparison cannot isolate how much of the gain came from the visual input itself.

The external studies tested whether the same finding held for different targets. In MSM, stacking ranked abnormal milling intervals less accurately than the logistic sensor model, with average precision across seeds of .757 and .829, respectively (Table~\ref{tab:msm-clean}). In ViFailback, the target was the outcome of a whole manipulation episode. Across three seeds, stacking yielded macro-AP of .621--.629 against .649 for logistic sensing. Initializing the visual encoder from REASSEMBLE gave mixed results across the matched runs (Tables~\ref{tab:vif-bothclass} and~\ref{tab:vif-matched}). In HANDOVER, averaging video and force--torque predictions on 151 trials yielded mean failure AP of .803, below the sensor component at .824 (Table~\ref{tab:handover-clean}). Participant-to-trial identifiers were unavailable, so the supplementary study could not construct the participant-disjoint predictions needed to train and evaluate its proposed stacker.

Taken together, the studies support a modular architecture in which each observation can be evaluated on its own before its predictions are combined with another stream. The value of that combination depended on the target, the available observations, and the strength of the sensor predictor. The results concern gear-condition estimation and retrospective classification of completed operations. Independent weld-quality detection remains the unanswered application of the original question.

\section{Supported Contributions}
The representation comparison (SQ1) favored sensor statistics on the main assembly task. In REASSEMBLE, the frozen ResNet18-vd backbone inside the pretrained RT-DETR-R18vd checkpoint supplied pooled frame features to a binary action classifier. Engineered sensor statistics outperformed the short-budget PatchTST time-series Transformer in the principal assembly comparison. Pooling features near the action boundary improved a separate exploratory visual result. The suitability of pretrained audio representations remains open because audio alignment did not support a predictive comparison.

For fusion choice (SQ2), the shared assembly cohort allowed feature concatenation, probability averaging, logistic stacking, and scalar gates conditioned on process context to be evaluated on the same actions. Feature concatenation reached AP of .412. The simpler prediction-level combinations ranked failures more accurately under the tested training procedures. The external results qualify that finding by identifying tasks on which even the learned combination fell below its sensor component.

For missing or degraded inputs (SQ3), a known missing stream could be handled by routing predictions to the available model. When a stream remained present but became unreliable, the two original gates evaluated under corruption showed no consistent advantage. Under strong sensor noise in REASSEMBLE, stacking fell to AP .222, below the intact visual predictor at .231. A subsequent gate that adjusted a static prediction using context improved clean performance in an exploratory comparison. Tests with matched corruption exposure left its robustness gain unresolved. This gate and the later visual pooling method were evaluated in separate experiments.

For evaluation beyond accuracy (SQ4), the distinction between ranking failures and estimating their probability affected model choice. On the seed-42 MSM comparison, stacking yielded a lower Brier score than logistic sensing (.056 and .062, respectively), although sensing ranked abnormal intervals more accurately. Grouped analyses, uncertainty estimates, a review of action labels, and component timings exposed different limits on reliability and practical use. The code and configuration records identify which branches and fusion rules produced these results. The literature matrix places them alongside prior welding, camera--sensor, and missing-input studies and narrows the open question to the performance of adaptive weighting under comparable training and calibration conditions.

\section{Principal Limitations}
Each dataset imposes a different evidential limit. PHM lacks an independent physical reference, and the preprocessing order in its separate four-arm diagnostic prevents a leakage-free interpretation of that experiment. The industrial and robotic datasets offer incomplete information about independence at the level of tools, objects, or participants. Complete-action and complete-episode predictions can use evidence that only becomes visible near the outcome. Deploying a warning system would require a prediction horizon that excludes such information.

The comparison also depends on how the models were fitted. Calibration, training budgets, and exposure to corruptions differed in some comparisons. Several follow-ups revisited previously evaluated populations, and the reported bootstrap intervals condition on the fitted models. The short-budget fine-tuning comparison did not resolve a mechanism for the observed differences. Some longer training runs and planned sensor-degradation evaluations remained incomplete. Where only run summaries were available, they limited the scope for independent checks of individual predictions.

\section{Priorities for Future Research}
An independent welding study is the most direct next test of the original question. It should specify the weld-level outcome and prediction time, pair the camera and process measurements to that decision, and reserve complete new welding sessions for assessment. An independent inspection should establish weld quality. For PHM, a physical measure of damage linked to a known tooth would likewise distinguish measurement validity from the ability to predict the photographic proxy. Independent review of ambiguous robotic outcomes would strengthen the corresponding classification targets.

Within each study, the informative comparison is between the strongest single-input predictor, equal averaging, a fitted static combination, and a gate trained with comparable calibration opportunities and corruption exposure. Repeated tests of known absence, unannounced corruption, temporal misalignment, and shifts between physical groups would show which failure conditions justify adaptive weighting. Prespecified margins for clean performance and grouped uncertainty estimates would make the cost of a robustness improvement visible.

The remaining transfer experiments require matched optimization budgets, normalization procedures, and held-out tasks. For operational monitoring, prospective tests should also measure warning time, false alarms, and end-to-end delay, and evaluate how operators act on uncertain predictions. These measurements would establish whether the modest ranking benefit observed on completed assembly actions translates into a useful decision while a weld is still in progress.
